\documentclass[12pt]{article}
\usepackage[margin=1in]{geometry}
\usepackage{enumitem}
\usepackage{titlesec}
\usepackage{parskip}
\usepackage{hyperref}
\usepackage{fontenc}

\title{\textbf{Carey 2009 Claude Guide}\\
\large Carey, John M. \textit{Legislative Voting and Accountability}.\\
Cambridge: Cambridge University Press, 2009.}
\author{}
\date{}

\begin{document}

\maketitle
\tableofcontents
\newpage

%--------------------------------------------------
\section{Central Claims}
%--------------------------------------------------

\begin{itemize}[leftmargin=*, itemsep=4pt]
    \item The book's animating question is about \textbf{accountability}: to whom are legislators actually responsive when they cast votes, and what determines whether they vote together with their party or defect from it? Carey treats the observable pattern of legislative voting---above all \textbf{party unity} on recorded (roll-call) votes---as the empirical footprint of underlying accountability relationships.
    \item The central theoretical contribution is the \textbf{theory of competing principals}. Building on principal--agent logic, Carey argues that legislators are agents who often answer to \emph{more than one} principal simultaneously---most importantly party leaders, but also voters/constituents, presidents, factional bosses, regional/state party organizations, and organized interests (unions, business, the Church). Party unity is high when a single principal dominates the legislator's incentives, and it breaks down when principals compete and pull the legislator in different directions.
    \item The key comparative-statics prediction: \textbf{party unity declines as the number and strength of competing principals increases}. Institutions that multiply or empower rival principals---candidate-centered electoral rules (open lists, personal-vote incentives, primaries), presidentialism, federalism/decentralized nomination, and bicameralism---systematically erode the party leadership's ability to secure unified voting.
    \item A second, distinctive contribution concerns the \textbf{transparency of voting}: whether votes are recorded and publicly visible is not neutral. Making votes transparent empowers whichever principal is in a position to \emph{observe and sanction} the legislator. Transparency therefore has no single, fixed effect on party unity---it strengthens party discipline when party leaders are the monitoring principal, but it can \emph{undermine} unity when transparency exposes legislators to a competing principal (e.g., constituents or interest groups) whose preferences diverge from the party's.
    \item This yields a pointedly normative implication that runs against a common reformist intuition: \textbf{more transparency is not unambiguously good for democratic accountability}. Because recorded voting reallocates influence toward whoever can watch, mandating transparency can enhance accountability to one principal at the expense of another, and can weaken the collective (party-based) accountability on which responsible-party models of democracy depend.
    \item Methodologically, the book is a large-scale \textbf{cross-national empirical study of roll-call voting}, assembling and analyzing legislative voting records from a wide range of democracies (with particular depth in Latin American presidential systems, alongside the U.S. and other cases) to measure party unity and test how institutional configurations of principals and transparency shape it.
    \item Carey positions the argument as a corrective to accounts that read party unity as evidence of shared preferences or ideology alone, or that treat it as a simple product of the parliamentary/presidential distinction. Unity, in his account, is an \emph{equilibrium outcome} of the structure of accountability---the configuration of principals and the visibility of the legislator's actions---rather than a fixed national or regime-type trait.
\end{itemize}

%--------------------------------------------------
\section{Key Concepts and Terms}
%--------------------------------------------------

\begin{itemize}[leftmargin=*, itemsep=4pt]
    \item \textbf{Accountability}: The condition of an agent being answerable to a principal who can observe the agent's behavior and reward or sanction accordingly. The book's organizing concept---legislative voting is analyzed as the behavior through which (and over which) accountability operates.
    \item \textbf{Principal--agent framework}: The workhorse logic---legislators are agents acting on behalf of principals who delegate authority and seek to control the agent's behavior via monitoring and sanctions (e.g., renomination, campaign resources, leadership posts, re-election).
    \item \textbf{Competing principals}: Carey's signature concept---the situation in which a legislator is accountable to two or more principals whose preferences over a vote diverge, forcing the legislator to choose whom to satisfy. The more numerous and evenly matched the competing principals, the lower expected party unity.
    \item \textbf{Party unity / party cohesion}: The degree to which members of the same party vote together on recorded votes. ``Unity'' is the observable co-voting pattern; ``cohesion'' often connotes shared underlying preferences. Carey stresses that observed unity can reflect discipline and incentive structure, not just shared preferences.
    \item \textbf{Unity scores (Rice index and related measures)}: Quantitative measures of within-party voting agreement computed from roll-call data. The \textbf{Rice index} (absolute difference between the shares of a party voting yea vs.\ nay on a given vote) is the classic building block; Carey computes and compares unity across parties, legislatures, and issue areas.
    \item \textbf{Transparency of voting}: Whether individual legislators' votes are recorded and made observable (to party leaders, the public, media, or interest groups). A central independent variable---transparency determines which principals can monitor and thus sanction the legislator.
    \item \textbf{Recorded / roll-call vote}: A vote in which each legislator's individual choice is registered and (typically) made public, as opposed to voice votes or secret ballots. The availability and selection of roll calls is both the book's data source and, substantively, a variable of interest.
    \item \textbf{Personal vote / candidate-centered incentives}: Electoral incentives (from Carey \& Shugart 1995) that push legislators to cultivate individual reputations distinct from the party label. Rules generating a strong personal vote (open-list PR, SNTV, primaries, intraparty competition) create a competing constituency principal and depress party unity.
    \item \textbf{Electoral rules / ballot structure}: Institutional features---list type (open vs.\ closed), district magnitude, pooling of votes, control over nominations and ballot rank---that determine whether party leaders or voters/rivals hold the decisive sanction over a legislator's career.
    \item \textbf{Nomination control}: Who decides whether a legislator is renominated (central party leaders vs.\ regional bosses vs.\ primary voters). Centralized nomination concentrates principal power in the leadership and raises unity; decentralized nomination fragments it.
    \item \textbf{Presidentialism as a competing principal}: In presidential systems the president---especially over co-partisans---can act as an additional principal (via patronage, agenda power, cabinet posts) rivaling or reinforcing legislative party leaders, complicating unity in ways distinct from parliamentary confidence-vote logic.
    \item \textbf{Federalism / decentralization}: The presence of subnational (state/provincial) party organizations and governors as principals with independent leverage over legislators' careers, cross-cutting national party leadership and lowering national party unity.
    \item \textbf{Agenda control}: Leaders' power to determine what reaches a vote and in what form; interacts with unity because leaders may avoid holding recorded votes that would reveal defection (selection of which votes become roll calls).
    \item \textbf{Discipline vs.\ cohesion (the observational-equivalence problem)}: The methodological challenge that high observed unity can result either from members genuinely agreeing (cohesion) or from leaders compelling agreement (discipline). Carey's institutional variation across principals is used to disentangle these.
    \item \textbf{Responsible party government}: The normative model in which cohesive, programmatic parties present clear alternatives so voters can hold the collective party accountable. Competing principals and certain transparency arrangements can undercut this model.
\end{itemize}

%--------------------------------------------------
\section{Core Cases and Data}
%--------------------------------------------------

\begin{itemize}[leftmargin=*, itemsep=4pt]
    \item \textbf{Cross-national roll-call dataset}: The empirical backbone is an assembled collection of legislative voting records spanning many legislatures, enabling comparison of party unity scores across institutional contexts. The comparative leverage comes from variation in electoral rules, presidential vs.\ parliamentary organization, federalism, and voting transparency.
    \item \textbf{Latin American presidential legislatures}: The region provides the richest variation and several of the book's key cases---e.g., \textbf{Argentina, Brazil, Bolivia, Chile, Mexico, Uruguay}, among others---because they combine presidentialism, diverse electoral rules (closed vs.\ open lists, varying district magnitude), and differing degrees of federalism and nomination centralization. Brazil and Argentina in particular illustrate how federalism and candidate-centered rules create competing principals that depress national party unity.
    \item \textbf{United States Congress}: A benchmark candidate-centered, weak-leadership-sanction system historically, used to illustrate how personal-vote incentives and decentralized nomination (primaries) coexist with (variable) party unity.
    \item \textbf{Parliamentary and other comparison cases}: Additional legislatures are drawn in to contrast high-unity, leadership-dominant settings (where the confidence relationship and centralized nomination make party leaders the dominant principal) with fragmented-principal settings.
    \item \textbf{Transparency ``natural'' variation}: Comparisons exploit differences (across and sometimes within legislatures over time) in whether and which votes are recorded/public, to examine how making votes transparent shifts unity depending on which principal gains the monitoring advantage.
    \item \textbf{Unit of analysis}: Primarily the party--vote and legislator--vote level, aggregated into party unity scores; institutional variables (electoral rule type, federalism, presidential powers, transparency regime) enter as explanators of variation in those scores.
\end{itemize}

%--------------------------------------------------
\section{Core Works Cited / Engaged}
%--------------------------------------------------

\begin{itemize}[leftmargin=*, itemsep=4pt]
    \item \textbf{John M. Carey and Matthew S. Shugart}, ``Incentives to Cultivate a Personal Vote: A Rank Ordering of Electoral Formulas'' (1995, \textit{Electoral Studies})---the foundational companion piece; its ranking of electoral rules by personal-vote incentives underpins the book's account of when constituents/rivals become competing principals.
    \item \textbf{Gary W. Cox and Mathew D. McCubbins}, \textit{Legislative Leviathan} (1993) and \textit{Setting the Agenda} (2005)---the leading party-cartel/agenda-control account of legislative parties, engaged as a partial and institutionally contingent story that competing-principals logic generalizes and qualifies.
    \item \textbf{Mathew D. McCubbins and others}, principal--agent and delegation theory in political science---the theoretical scaffolding Carey extends from single-principal to multiple-principal settings.
    \item \textbf{Stuart Rice}, foundational work on measuring legislative voting agreement (the \textbf{Rice index})---the classic measurement tool for party unity that Carey adopts and refines.
    \item \textbf{Keith Poole and Howard Rosenthal}, spatial/ideal-point (NOMINATE) scaling of roll-call votes---the dominant methodology for analyzing recorded votes, part of the methodological backdrop against which Carey's institutional, accountability-centered reading of roll calls is positioned.
    \item \textbf{Scott Morgenstern}, \textit{Patterns of Legislative Politics} (2004) and related comparative work on Latin American and U.S.\ legislative voting---a close interlocutor on cross-national party unity and its determinants.
    \item \textbf{Gary Cox}, \textit{The Efficient Secret} (1987)---on the historical development of party discipline and the role of vote visibility/procedure, relevant to the transparency argument.
    \item \textbf{Barbara Geddes, Scott Mainwaring, Matthew Shugart}, and the broader comparative literature on presidentialism, parties, and electoral systems in Latin America---the empirical and institutional literature the book's cases are embedded in.
    \item \textbf{Responsible-party-government tradition} (e.g., the APSA 1950 report and its descendants)---the normative benchmark against which Carey assesses the accountability consequences of competing principals and transparency.
\end{itemize}

%--------------------------------------------------
\section{Part I: The Accountability Problem and the Theory of Competing Principals}
%--------------------------------------------------

\begin{itemize}[leftmargin=*, itemsep=4pt]
    \item \textbf{Framing}: Opens by reframing legislative voting as a window onto accountability. Rather than asking simply ``are parties unified?'', Carey asks ``to whom is the individual legislator accountable, and how does that shape the vote?''---shifting the object of explanation from regime type to the structure of principal--agent relations.
    \item \textbf{From single principal to competing principals}: Standard principal--agent treatments of legislators assume (implicitly) one dominant principal---typically the party leadership or the median voter. Carey's move is to take seriously that legislators routinely face \emph{several} principals at once, and that the interesting variation in legislative behavior comes from the \emph{configuration} of those principals.
    \item \textbf{The core prediction stated}: When principals are aligned (they want the same vote) or when one clearly dominates (holds the decisive career sanction), the legislator votes predictably and, if the dominant principal is the party, unity is high. When principals compete---divergent preferences, comparably strong sanctions---unity falls, defection rises, and behavior becomes harder to predict from party label alone.
    \item \textbf{Sources of principal power}: Control over renomination, position on the ballot/list, campaign resources, leadership and committee assignments, cabinet posts and patronage (in presidential systems), and ultimately re-election. Whoever controls the sanctions that matter most for a legislator's career is the principal whose preferences the legislator will privilege.
    \item \textbf{Why this generalizes existing theories}: Agenda-cartel and party-discipline accounts are recovered as the special case in which the party leadership is the dominant principal; median-voter/electoral-connection accounts are the special case in which constituents dominate. Competing-principals theory nests both and predicts the intermediate, mixed, and unstable cases.
\end{itemize}

%--------------------------------------------------
\section{Part II: Measuring Party Unity and Its Institutional Determinants}
%--------------------------------------------------

\begin{itemize}[leftmargin=*, itemsep=4pt]
    \item \textbf{Measurement}: Develops and applies unity scores (Rice-index-based and related measures) computed from roll-call records, attentive to the pitfalls of comparing across legislatures with different agendas, different sets of recorded votes, and different party sizes.
    \item \textbf{The discipline-vs.-cohesion problem}: Confronts the observational equivalence between members voting together because they agree and members voting together because they are compelled. The empirical strategy leans on \emph{institutional variation in principals}: if unity tracks the configuration of principals (rather than only issue content or ideology), that supports the accountability-based, discipline-inclusive interpretation.
    \item \textbf{Electoral rules and the personal vote}: Draws on the Carey--Shugart (1995) ranking to predict that candidate-centered rules (open lists, intraparty competition, primaries, large personal-vote incentives) empower constituents and co-partisan rivals as competing principals, lowering unity; party-centered rules (closed lists, centralized nomination) concentrate power in leadership and raise unity.
    \item \textbf{Presidentialism}: Examines how a president---particularly vis-\`a-vis co-partisans---functions as an additional principal through patronage, cabinet appointments, and agenda/budget powers, sometimes reinforcing and sometimes cross-cutting legislative party leaders' control over the same legislators.
    \item \textbf{Federalism and decentralized nomination}: Shows how governors, state parties, and regional bosses hold independent leverage over legislators' careers, fragmenting the principal structure and systematically depressing \emph{national} party unity (with Brazil and Argentina as leading illustrations).
    \item \textbf{Cross-national findings}: The comparative pattern is that unity is highest where institutions concentrate the decisive career sanction in a single (party-leadership) principal, and lowest where electoral rules, presidential patronage, and federal decentralization multiply competing principals.
\end{itemize}

%--------------------------------------------------
\section{Part III: The Transparency of Voting}
%--------------------------------------------------

\begin{itemize}[leftmargin=*, itemsep=4pt]
    \item \textbf{Transparency as a variable, not a constant good}: The book's most provocative move is to treat the visibility of votes as an institutional choice with distributional consequences. Recording and publicizing votes reallocates influence toward whichever principal can \emph{observe} the legislator and act on what it sees.
    \item \textbf{Conditional effect on unity}: Transparency does not have a single sign. Where the party leadership is the monitoring principal, recorded voting \emph{tightens} discipline (leaders can detect and punish defectors), raising unity. Where transparency instead exposes the legislator to a competing principal---constituents, media, organized interests---whose preferences diverge from the party's, recorded voting can \emph{lower} unity by activating cross-pressure.
    \item \textbf{Selection of roll calls}: Because leaders often influence which votes are recorded, observed roll-call unity is partly endogenous---leaders may avoid recorded votes that would publicize costly defection, so the set of roll calls is not a random sample of legislative decisions. Carey treats this selection as substantively meaningful evidence about accountability, not merely a nuisance.
    \item \textbf{Normative payoff}: Reform advocates typically assume that more recorded, public voting straightforwardly improves accountability. Carey argues the effect depends on \emph{whose} accountability---transparency can strengthen accountability to constituents while weakening the collective, programmatic accountability of parties (responsible-party government), or vice versa. There is a genuine trade-off, not a free lunch.
    \item \textbf{Implication for institutional design}: Choices about vote transparency should be understood as choices about which principal to empower. Designers who value cohesive programmatic parties and those who value individual constituency responsiveness will (rationally) want different transparency regimes.
\end{itemize}

%--------------------------------------------------
\section{Conclusion and Contributions}
%--------------------------------------------------

\begin{itemize}[leftmargin=*, itemsep=4pt]
    \item \textbf{Synthesis}: Party unity in legislative voting is best understood as an equilibrium outcome of the structure of accountability---the number, strength, and alignment of the principals to whom a legislator answers, and the visibility of the legislator's votes to those principals---rather than as a fixed feature of a country, party, or regime type.
    \item \textbf{Theoretical contribution}: Generalizes principal--agent accounts of legislative behavior from a single dominant principal to \textbf{competing principals}, nesting agenda-cartel/party-discipline and electoral-connection theories as special cases and explaining the intermediate and unstable cases they miss.
    \item \textbf{Empirical contribution}: Provides one of the most systematic cross-national analyses of roll-call-based party unity, especially valuable for its coverage of Latin American presidential systems, and links variation in unity to concrete institutional features (list type, nomination control, presidentialism, federalism, transparency).
    \item \textbf{Normative contribution}: Complicates the reformist presumption in favor of transparency by showing that recorded voting redistributes accountability rather than simply increasing it---an important caution for institutional designers and a bridge to debates over responsible party government.
    \item \textbf{Legacy}: Consolidates ``competing principals'' as a standard analytical lens in comparative legislative studies and in the study of party unity, agenda control, and the accountability consequences of electoral and legislative institutions.
\end{itemize}

\end{document}
